\subsection{Introduction}
With the basic formalism and notation for continuous variable states established, we now address the problem of discriminating between such states, focusing particularly on Gaussian states.\\

A continuous variable state is classified as Gaussian if its Wigner function -- or equivalently, its characteristic function -- takes a Gaussian form. Mathematically, this is expressed as:

\begin{equation}
    W[\hat{\rho}](\vec{R}) = \frac{1}{(2\pi)^n \sqrt{\det \sigma}} \exp\left(-\frac{1}{2} (\vec{R} - \bar{R})^T \sigma^{-1} (\vec{R} - \bar{R})\right)
\end{equation}

where,
\begin{itemize}
    \item $W[\hat{\rho}](\vec{R})$: The Wigner function tells us the `probability-like' value at point $\vec{R}$ in phase space.
    \item $\vec{R}$: This is a vector containing the positions and momenta of all $n$ modes. For example, for one mode, $\vec{R} = (q, p)$.
    \item $\bar{R}$: The average (mean) value of $\vec{R}$. This is the center of the bell curve.
    \item $\sigma$: The covariance matrix. It tells us how spread out the state is in different directions.
    \item $\det \sigma$: The determinant of the covariance matrix. It measures the overall volume of the distribution.
    \item $n$: The number of modes (like the number of particles or light beams).
\end{itemize}

Alternatively,

\begin{equation}
    \chi[\hat{\rho}](\vec{\Lambda}) = \exp\left(-\frac{1}{2} \vec{\Lambda}^T \sigma \vec{\Lambda} + i \vec{\Lambda}^T \bar{R}\right)
\end{equation}

where,
\begin{itemize}
    \item $\chi[\hat{\rho}](\vec{\Lambda})$: The characteristic function. It is like the `moment-generating' function in statistics.
    \item $\vec{\Lambda}$: A vector in phase space that is conjugate to $\vec{R}$ (like frequency is conjugate to time).
    \item $\sigma$: The same covariance matrix from Equation (3.1).
    \item $\bar{R}$: The same mean vector from Equation (3.1).
\end{itemize}

Thus, a Gaussian state $\hat{\rho}$ is completely characterized by its covariance matrix $\sigma$ and its mean vector $\bar{R}$. The covariance matrix gains further physical significance through Williamson's Theorem, which decomposes it into a symplectic matrix and a diagonal matrix of symplectic eigenvalues. Importantly, all entanglement properties of a Gaussian state are contained within $\sigma$. This is because local operations do not affect entanglement, so the mean vector $\bar{R}$ can be arbitrarily adjusted through phase space translations without altering the entanglement of $\hat{\rho}$. Conversely, generating entanglement necessarily requires a nonlocal operation on $\hat{\rho}$.

\subsection{Normal form decomposition}
A useful theorem for expressing the covariance matrix of a Gaussian state is Williamson's Theorem, stated as follows:

\begin{theorem}[Williamson's Theorem]
    Given a real $2n \times 2n$ symmetric and positive matrix $\sigma$, there exists a real and symplectic $2n \times 2n$ matrix $S$ and a real $n \times n$ diagonal and positive matrix $W$ such that
    \begin{equation}
        \sigma = S^T W S
    \end{equation}
    where
    \begin{equation}
        W = \bigoplus_{k=1}^{n} d_k \mathbb{I}_2
    \end{equation}
    and
    \begin{equation}
        \mathbb{I}_2 = \begin{pmatrix} 1 & 0 \\ 0 & 1 \end{pmatrix}
    \end{equation}
\end{theorem}

\begin{proof}
The proof proceeds in several steps. We will construct the symplectic matrix $S$ that diagonalizes $\sigma$ in the symplectic sense.\\

Since $\sigma$ is real, symmetric, and positive-definite, it can be diagonalized by an orthogonal matrix $O \in O(2n)$:
\begin{equation}
    \sigma = O^T D O,
\end{equation}
where $D = \operatorname{diag}(\lambda_1, \lambda_2, \ldots, \lambda_{2n})$ with $\lambda_i > 0$. However, this is an orthogonal diagonalization, not a symplectic one. Our goal is to find a symplectic matrix $S$ that brings $\sigma$ to the block-diagonal form $W$.\\

Define the matrix
\begin{equation}
    M = i \Omega \sigma,
\end{equation}
where $\Omega$ is the symplectic form:
\begin{equation}
    \Omega = \bigoplus_{k=1}^{n} \omega, \qquad \omega = \begin{pmatrix} 0 & 1 \\ -1 & 0 \end{pmatrix}.
\end{equation}

Consider the matrix $M = i\Omega\sigma$. Its transpose is:
\begin{equation}
    M^T = (i\Omega\sigma)^T = i \sigma^T \Omega^T = i \sigma (-\Omega) = -i \sigma \Omega.
\end{equation}

Now note that:
\begin{equation}
    \Omega^{-1} M \Omega = \Omega^{-1} (i \Omega \sigma) \Omega = i \sigma \Omega.
\end{equation}

Since $\sigma \Omega$ is similar to $-\Omega \sigma$ (because $\Omega^T = -\Omega$), the eigenvalues of $M$ come in $\pm$ pairs. With the factor $i$, the eigenvalues are of the form $\pm i d_k$ where $d_k > 0$.\\

Let the eigenvalues of $M$ be $\pm i d_1, \pm i d_2, \ldots, \pm i d_n$ with $d_k > 0$. For each eigenvalue $i d_k$, there exists a complex eigenvector $u_k$ such that:
\begin{equation}
    i \Omega \sigma u_k = i d_k u_k \quad \Rightarrow \quad \Omega \sigma u_k = d_k u_k.
\end{equation}

Similarly, for eigenvalue $-i d_k$, we have an eigenvector $v_k$ satisfying $\Omega \sigma v_k = -d_k v_k$.\\

Using the fact that $\sigma$ is symmetric and $\Omega$ is skew-symmetric, one can show that $v_k$ is proportional to the complex conjugate of $u_k$. We choose them such that:
\begin{equation}
    v_k = i \sigma u_k.
\end{equation}
This choice ensures the correct normalization.\\

We normalize the eigenvectors such that:
\begin{equation}
    u_k^\dagger \sigma u_k = 1.
\end{equation}

Then, using properties of $\Omega$ and $\sigma$, we also get:
\begin{equation}
    u_k^\dagger \Omega u_k = i.
\end{equation}

Define the real vectors $x_k$ and $y_k$ for $k = 1, \ldots, n$ as:
\begin{equation}
    x_k = \sqrt{2} \, \operatorname{Re}(u_k), \qquad y_k = \sqrt{2} \, \operatorname{Im}(u_k).
\end{equation}

One can verify that these vectors satisfy:
\begin{align}
    x_k^T \sigma x_k &= y_k^T \sigma y_k = d_k, \\
    x_k^T \sigma y_k &= 0,
\end{align}
and
\begin{align}
    x_k^T \Omega x_k &= y_k^T \Omega y_k = 0, \\
    x_k^T \Omega y_k &= 1.
\end{align}

Let $S$ be the $2n \times 2n$ matrix whose columns are, in order:
\begin{equation}
    S = \begin{pmatrix} x_1 & y_1 & x_2 & y_2 & \cdots & x_n & y_n \end{pmatrix}.
\end{equation}

From the properties above:
\begin{itemize}
    \item Symplectic condition: $S^T \Omega S = \Omega$ (since $x_k^T \Omega y_k = 1$ and all other pairings are zero).
    \item Diagonalization of $\sigma$: $S^T \sigma S = \bigoplus_{k=1}^n d_k \mathbb{I}_2 = W$.
\end{itemize}

Thus, we have constructed a symplectic matrix $S$ such that:
\begin{equation}
    S^T \sigma S = W \quad \Rightarrow \quad \sigma = (S^{-1})^T W S^{-1}.
\end{equation}

Since $S$ is symplectic, $S^{-1}$ is also symplectic. Renaming $S^{-1}$ as $S$ (or absorbing the inverse into the definition), we obtain:
\begin{equation}
    \sigma = S^T W S
\end{equation}
with $S \in Sp(2n, \mathbb{R})$ and $W$ diagonal as required. This completes the proof.
\end{proof}


The elements $d_k$ of the matrix $W$ are called the symplectic eigenvalues and can be calculated from the spectrum of $i\Omega\sigma$, while the matrix $S$ is the symplectic matrix that performs the symplectic decomposition. Finding the symplectic eigenvalues of a covariance matrix is usually straightforward, but determining the symplectic matrix that performs the decomposition is often difficult. Performing the symplectic decomposition of Gaussian states is a key part of the process of determining the distinguishability of those states. For example, it turns out that in order to calculate the quantum Chernoff bound for Gaussian states -- an asymptotic distinguishability measure for states that will be used extensively -- we must find the symplectic decomposition of the states involved.\\

The phenomenological statement implied by equation (3.3) is that every Gaussian state $\hat{\rho}_G$ can be obtained from a thermal state $\hat{\rho}_{\text{th}}$ (described by a diagonal covariance matrix as in equation (2.96)) by performing the unitary transformation $U_S$ associated with the symplectic matrix $S$ that performs the symplectic decomposition:

\begin{equation}
    \hat{\rho}_G = U_S \hat{\rho}_{\text{th}} U_S^\dagger
\end{equation}

where $\hat{\rho}_{\text{th}} = \hat{\rho}_{\text{th},1} \otimes \hat{\rho}_{\text{th},2} \otimes \cdots \otimes \hat{\rho}_{\text{th},n}$ is a product of thermal states $\hat{\rho}_{\text{th},k}$ for $k = 1, \ldots, n$.

Now, we derive the above equation.\\

From Williamson's Theorem, we have:

\begin{equation}
    \sigma = S^T W S \tag{3.3}
\end{equation}

where:
\begin{itemize}
    \item $\sigma$ is the covariance matrix of an arbitrary Gaussian state $\hat{\rho}_G$
    \item $S$ is a symplectic matrix ($S \in Sp(2n, \mathbb{R})$)
    \item $W$ is a diagonal matrix of the form:
    \begin{equation}
        W = \bigoplus_{k=1}^{n} d_k \mathbb{I}_2, \qquad \mathbb{I}_2 = \begin{pmatrix} 1 & 0 \\ 0 & 1 \end{pmatrix}
    \end{equation}
\end{itemize}

For a single-mode thermal state $\hat{\rho}_{\text{th},k}$ with average photon number $N_k$, the covariance matrix is (from equation 2.96):

\begin{equation}
    \sigma_{\text{th},k} = \frac{1}{4} \begin{pmatrix} 2N_k + 1 & 0 \\ 0 & 2N_k + 1 \end{pmatrix} = \left(N_k + \frac{1}{2}\right) \mathbb{I}_2
\end{equation}

Comparing this with the block form of $W$, we identify:

\begin{equation}
    d_k = N_k + \frac{1}{2}
\end{equation}

Thus, the matrix $W$ can be written as:

\begin{equation}
    W = \bigoplus_{k=1}^{n} \left(N_k + \frac{1}{2}\right) \mathbb{I}_2
\end{equation}

This is exactly the covariance matrix of a \textbf{product of thermal states}:

\begin{equation}
    \hat{\rho}_{\text{th}} = \hat{\rho}_{\text{th},1} \otimes \hat{\rho}_{\text{th},2} \otimes \cdots \otimes \hat{\rho}_{\text{th},n}
\end{equation}

whose covariance matrix is:

\begin{equation}
    \sigma_{\text{th}} = W
\end{equation}

Now, we want to find a unitary transformation $U_S$ that maps the thermal state $\hat{\rho}_{\text{th}}$ to the target Gaussian state $\hat{\rho}_G$.

In phase space, a unitary transformation generated by a quadratic Hamiltonian corresponds to a symplectic transformation on the covariance matrix (see Appendix C of the thesis):

\begin{equation}
    \sigma \rightarrow S \sigma S^T
\end{equation}

More precisely, if $\hat{\rho}' = U_S \hat{\rho} U_S^\dagger$, then the covariance matrix transforms as:

\begin{equation}
    \sigma' = S \sigma S^T
\end{equation}

where $S$ is the symplectic matrix associated with the unitary $U_S$.\\

Let us apply the symplectic transformation $S$ to the thermal state $\hat{\rho}_{\text{th}}$:

\begin{equation}
    \hat{\rho}_G = U_S \hat{\rho}_{\text{th}} U_S^\dagger
\end{equation}

The covariance matrix of $\hat{\rho}_G$ becomes:

\begin{equation}
    \sigma_G = S \sigma_{\text{th}} S^T
\end{equation}

But $\sigma_{\text{th}} = W$, so:

\begin{equation}
    \sigma_G = S W S^T
\end{equation}

Williamson's Theorem states:

\begin{equation}
    \sigma = S^T W S
\end{equation}

Notice the difference: Williamson's theorem has $S^T W S$, while we have $S W S^T$. However, since $S$ is symplectic, its inverse is given by:

\begin{equation}
    S^{-1} = \Omega S^T \Omega^{-1}
\end{equation}

More simply, if $S$ is symplectic, then $S^T$ is also symplectic. Therefore, we can rename $S^T$ as a new symplectic matrix $S'$. This means:

\begin{equation}
    S W S^T = (S^T)^T W (S^T) = S'^T W S'
\end{equation}

Thus, both forms are equivalent. Williamson's theorem guarantees the existence of some symplectic matrix such that $\sigma = S^T W S$. Therefore, there exists a unitary transformation $U_S$ such that:

\begin{equation}
    \hat{\rho}_G = U_S \hat{\rho}_{\text{th}} U_S^\dagger
\end{equation}

Thus, we arrive at the equation:

\begin{equation}
    \hat{\rho}_G = U_S \hat{\rho}_{\text{th}} U_S^\dagger
\end{equation}

where:
\begin{itemize}
    \item $\hat{\rho}_G$ is any Gaussian state
    \item $\hat{\rho}_{\text{th}} = \hat{\rho}_{\text{th},1} \otimes \cdots \otimes \hat{\rho}_{\text{th},n}$ is a product of thermal states
    \item $U_S$ is the unitary transformation associated with the symplectic matrix $S$
\end{itemize}
